\documentclass[11pt]{article}
\usepackage[margin=25mm]{geometry}
\usepackage{fontspec}
\setmainfont{Latin Modern Roman}
\usepackage{amsmath,amssymb,mathtools}
\usepackage{unicode-math}
\setmathfont{Latin Modern Math}
\usepackage{xurl,hyperref}
\hypersetup{colorlinks=true,urlcolor=blue}
\setcounter{secnumdepth}{0}
\setlength{\parindent}{0pt}
\setlength{\parskip}{5pt}
\setlength{\emergencystretch}{3em}
\newcommand{\sref}[2]{\hyperlink{source:#1:#2}{[S#2]}}
\newcommand{\cataloguescope}[1]{\paragraph{Recorded target.}#1}
\begin{document}
\section{Upgrading a completary multifuncoid}\label{statement}

\subsection{Definitions and mathematical statement}

\paragraph{Shared notation.}$\mathbb N=\{1,2,\ldots\}$, and $\mathbb Z,\mathbb Q,\mathbb R,\mathbb C$ denote the integers, rationals, reals and complex numbers. Any other conventions are specified locally.


For a set $X$, let $\mathcal F(X)$ be all filters on $\mathcal P(X)$, including the improper filter $\bot_X=\mathcal P(X)$. A filter is a nonempty family closed under finite intersections and supersets. Order filters by $\mathcal A\sqsubseteq\mathcal B$ when $\mathcal A\supseteq\mathcal B$ as families. Thus $\bigsqcup$ is intersection of filter families and $\bigsqcap$ is the filter generated by their union; the empty join is $\bot_X$ and the empty meet is $\top_X=\{X\}$. For $E\subseteq X$, $\uparrow E=\{U\subseteq X:E\subseteq U\}$ is its principal filter. Write $\mathcal A\mathrel{\#}\mathcal B$ when every $U\in\mathcal A$ and $V\in\mathcal B$ have $U\cap V\ne\varnothing$, equivalently $\mathcal A\sqcap\mathcal B\ne\bot_X$.

For an index set $I$ and join-semilattices $L_i$ with bottom $0_i$, their product is ordered coordinatewise. Call $F\subseteq\prod_{i\in I}L_i$ completary when, for every $u,v$ in that product,
\[
 (u_i\vee v_i)_{i\in I}\in F
 \quad\Longleftrightarrow\quad
 \exists c:I\to\{0,1\}:\ (w_i)_{i\in I}\in F,
 \quad w_i=\begin{cases}u_i&c(i)=0,\\v_i&c(i)=1,\end{cases}
\]
and no tuple having a bottom coordinate belongs to $F$. This is the source's definition of a completary multifuncoid.Fix any set $X$ and any index set $I$. Let $\mathcal P_f(X)=\{\uparrow A:A\subseteq X\}\subseteq\mathcal F(X)$ be the principal filters, with inherited order. For $u\in\mathcal F(X)^I$, define
\[
 \operatorname{up}(u)=\{v\in\mathcal P_f(X)^I:u_i\sqsubseteq v_i\text{ for all }i\in I\},
 \quad E^*F=\{u\in\mathcal F(X)^I:\operatorname{up}(u)\subseteq F\}.
\]
Here $\operatorname{up}$ and $E^*$ refer to the fixed filtrator $(\mathcal F(X)^I,\mathcal P_f(X)^I)$.
\paragraph{Statement recorded in the supplied source.}
For every $X,I$ and every completary multifuncoid $F\subseteq\mathcal P_f(X)^I$, is its upgraded set $E^*F\subseteq\mathcal F(X)^I$ completary for the filter-lattice joins? The quantifier includes infinite index sets. \sref{statement}{2}
\subsection{Short English statement}
Does upgrading a completary multifuncoid from principal filters to all filters preserve its defining simultaneous join rule, even with infinitely many coordinates?
\subsection{Sources}
\begin{description}
\item[\hypertarget{source:statement:1}{S1}] ULAM.AI and the UnsolvedMath Contributors, UnsolvedMath: A Curated Collection of Open Mathematics Problems, catalogue entry. CC BY 4.0. \url{https://huggingface.co/datasets/ulamai/UnsolvedMath}.
\item[\hypertarget{source:statement:2}{S2}] Open Problem Garden, Upgrading a completary multifuncoid, catalogue entry. Original problem statement, full mathematical discussion, bibliography and comments. \url{https://www.openproblemgarden.org/op/upgrading_a_completary_multifuncoid}.
\end{description}
\label{end:statement}
\end{document}
